\section{Conclusions}

In conclusion, none of the three LLMs reliably reproduced human CLAIM assessments across the 60 evaluated papers and no single model consistently achieved the highest agreement with human assessments. ChatGPT generally showed the largest discrepancies, whereas Claude and Gemini performed better in different domains. Correlations with human scores were low and non-significant, indicating that the models also rank papers differently from humans.

Within the scope of this study, general-purpose LLMs are not yet reliable substitutes for human CLAIM assessment under the conditions tested, and their use in peer review should remain supervised. Item-level analyses are needed to identify where their judgments can be trusted.

\subsection{Limitations and improvements}
This study has several limitations. First, the lack of item-level human assessments prevented us from identifying which specific CLAIM items were responsible for disagreements. Moreover, human scores may be subject to some degree of subjectivity, meaning that our results reflect disagreement rather than absolute accuracy. Second, the sample was limited to the first 60 available papers from a single imaging field and used CLAIM 2020, limiting the generalization of the findings. Finally, only one prompt and one run per paper were tested using the free versions of the models, preventing assessment of prompt sensitivity and run-to-run variability.

Future work should focus on item-level comparisons, include human-human agreement as a reference, and evaluate multiple prompts and repeated runs. Larger datasets, additional imaging domains, CLAIM 2024 and models with extended reasoning should also be explored. Finally, simple baselines and calibration or fine-tuning strategies could help determine whether LLM reliability can be improved and establish their added value in CLAIM assessment.
